\section{Summation by Parts}
\begin{theorem}[Summation by Parts]
\label{th:somma_per_parti}%
Let $a_k$ and $B_k$ be sequences (real or complex).
Define
\[
  A_n = \sum_{k=0}^{n-1} a_k, \qquad
  b_n = B_{n+1} - B_n
\]
then
\mymargin{summation by parts}%
\index{summation!by parts}%
\begin{equation}\label{eq:somma_per_parti}
 \sum_{k=m}^{n-1} a_k B_k = A_n B_n - A_m B_m - \sum_{k=m}^{n-1} A_{k+1}b_k.
\end{equation}
\end{theorem}
%
\begin{proof}
Observe that $a_n = A_{n+1}-A_n$, thus
\begin{align*}
  A_n B_n - A_m B_m
  &= \sum_{k=m}^{n-1} (A_{k+1}B_{k+1}-A_k B_k)\\
  &= \sum_{k=m}^{n-1} (A_{k+1}B_{k+1}-A_{k+1}B_k + A_{k+1}B_k - A_k B_k)\\
  &= \sum_{k=m}^{n-1} A_{k+1}b_k + \sum_{k=m}^{n-1} a_k B_k.
\end{align*}
\end{proof}

If we take $a_k=(-1)^k$, it can be observed that $A_n = \sum_{k=0}^{n-1} (-1)^k$
is a bounded sequence since $A_n = 1$ if $n$ is odd while
$A_n=0$ if $n$ is even.
If instead we choose a sequence $B_n$ that is positive,
decreasing, and infinitesimal
and set $b_n = B_{n+1}-B_n$
we have
\[
  \sum_{k=0}^{+\infty} \abs{b_k}
  = \lim_{n\to+\infty}\sum_{k=0}^{n-1} (B_k-B_{k+1})
  = \lim_{n\to +\infty} (B_0 - B_n) = B_0 < +\infty.
\]
Thus, the following theorem is an extension of the
Leibniz criterion for alternating series.

\begin{theorem}[Dirichlet's Criterion]%
\label{th:dirichlet}%
\index{criterion!of Dirichlet}%
\index{Dirichlet!criterion of}%
Let $a_n$ and $B_n$ be sequences (real or complex)
and set $\displaystyle A_n = \sum_{k=0}^{n-1} a_k$,
$b_n = B_{n+1} - B_n$.
If $A_n$ is bounded,
$B_n\to 0$
and $\sum \abs{b_n} < +\infty$
then the series $\sum a_k B_k$ is convergent.
\end{theorem}
%
\begin{proof}
By the summation by parts formula~\eqref{eq:somma_per_parti}
 we have
\begin{equation}\label{eq:3498954}
 \sum_{k=0}^{n-1} a_k B_k
 = A_n B_n - A_0 B_0 - \sum_{k=0}^{n-1} A_{k+1}b_k.
\end{equation}
The first term $A_n B_n$ tends to zero as $n\to +\infty$
since it is the product of a bounded sequence and an infinitesimal.
The second term is constant.
The series $\sum A_{k+1} b_k$ is absolutely convergent
since $\abs{A_n}\le L$ is bounded
and $\sum b_n$ is absolutely convergent, we have
\[
  \sum \abs{A_{k+1}} \cdot \abs{b_k}
  \le L \cdot \sum \abs{b_k} < +\infty.
\]

Thus, the limit of the sum on the right-hand side converges, and therefore
the sum on the left-hand side
of equation~\eqref{eq:3498954}
is convergent as $n\to +\infty$.
\end{proof}

\begin{example}\label{ex:serie_log}
Given $z\in \CC$, $\abs{z} \le 1$, $z\neq 1$, the series
\[
  \sum_{k=1}^{+\infty} \frac{z^k}{k}
\]
is convergent.
\end{example}
\begin{proof}
Note that for $\abs{z}<1$, one can easily apply
the ratio or root test to the series of moduli.
But for $\abs{z}=1$, those tests do not apply, and
one must instead use Theorem~\ref{th:dirichlet}.

Setting $a_k=z^k$, it is observed that $\sum a_k$ is a
geometric series and (since $z\neq 1$) we have
\[
  A_n = \sum_{k=0}^{n-1} z^k = \frac{1-z^n}{1-z}
\]
which is bounded, indeed:
\[
  \abs{A_n} = \frac{\abs{1-z^n}}{\abs{1-z}} \le \frac{1+\abs{z^n}}{\abs{1-z}}
  = \frac{2}{\abs{1-z}}.
\]
While setting $B_n = 1/n$, it is clear that, being $B_n$ real,
decreasing, we have
\begin{align*}
\sum_{k=1}^{+\infty} \abs{B_{k+1}-B_{k}}
&= \lim_{n\to +\infty}\sum_{k=1}^{n-1} (B_k-B_{k+1})
= \lim_{n\to +\infty}(B_1 - B_n)\\
&= \lim_{n\to +\infty}\enclose{1 - \frac{1}{n}} = 1
< +\infty.
\end{align*}
Thus, Theorem~\ref{th:dirichlet} is applied
to obtain the convergence
of the given series.

Note that if $\abs{z}>1$, the series in question does not
converge because the general term $z^k/k$ is not infinitesimal.
For $z=1$, the harmonic series is obtained, which also does not converge.
\end{proof}

\begin{exercise}
  The series
  \[
    \sum_{k=1}^{+\infty} \frac{\sin k}{k}
  \]
  is convergent.
  \end{exercise}
  \begin{proof}
  We apply Theorem~\ref{th:dirichlet}.
  Setting $a_k = \sin k$ and $B_k=1/k$
  we have
  \[
    A_n = \sum_{k=0}^{n-1} \sin k = \Im \sum_{k=0}^{n-1} e^{ik}.
  \]
  We then observe that $e^{ik}=(e^i)^k$, and thus $A_n$ is the imaginary part
  of the sum of a geometric series. 
  It can therefore be calculated explicitly
  \[
    A_n = \Im \frac{1-(e^i)^n}{1-e^i}
  \]
  from which
  \[
    \abs{A_n} \le \abs{\frac{1-e^{in}}{1-e^i}} \le \frac{1+\abs{e^{in}}}{\abs{1-e^i}}
    = \frac{2}{\abs{1-e^i}}
  \]
  and thus $A_n$ is bounded.
  
  On the other hand, setting $B_k = 1/k$, it is clear that $B_k$ is
  decreasing and infinitesimal.
  \end{proof}